\chapter{Estado del arte}
\label{cap:estado-arte}

Nuestro problema a tratar vive en el cruce de tres líneas que han avanzado por separado: los agentes LLM que ya operan en los mercados, la idea de supervisar un sistema mientras decide, y los modelos clásicos que describen el régimen de mercado y su riesgo. Recorremos las tres y dejamos ver el hueco que STRATA ocupa.

\section{Agentes LLM en finanzas}
\label{sec:agentes-llm}

Un gran modelo de lenguaje (LLM) es una red que predice la siguiente palabra a partir del contexto, mediante un mecanismo de atención \cite{vaswani2017attention}. Por sí solo solo escribe: se vuelve un \emph{agente} cuando lo envolvemos en un bucle que razona, llama a herramientas y actúa sobre un entorno \cite{yao2023react, wang2024survey}. De ahí salieron sistemas que deciden inversiones en lugar de redactar \cite{yang2023fingpt}. El foco se desplazó pronto de si un modelo entiende los mercados a cómo coordinar varios, y hoy domina la arquitectura multiagente: varias personalidades inversoras deliberan y un gestor modera la posición final \cite{xiao2024tradingagents, zhao2025alphaagents}. El control de riesgo que incorporan vive dentro del propio sistema y comparte sus sesgos: el equipo de riesgo de TradingAgents razona en lenguaje natural \cite{xiao2024tradingagents} y AlphaAgents ajusta la cartera a la tolerancia del inversor \cite{zhao2025alphaagents}, pero ninguno audita con un modelo estadístico externo si cada decisión es coherente con el mercado observado.

 Los agentes LLM ejecutan las instrucciones con fidelidad antes que maximizar beneficio, aun cuando eso les hace perder \cite{lopezlira2025cantrade}, y su confianza declarada está peor calibrada que la de un humano: se muestran sobreconfiados justo cuando fallan \cite{sung2025grace, damani2025rlcr}. Cambian de estilo inversor sin que el cambio responda a su propio criterio \cite{li2026behavioral} y resultan manipulables ante datos diseñados para engañarlos \cite{yan2025tradetrap}. Un agente rígido cuya certeza no es fiable pide un supervisor exógeno. Y la evaluación es frágil: muchos resultados optimistas no aguantan una validación temporal honesta \cite{li2025deepfund}, y el \emph{look-ahead} infla las métricas cuando el periodo de prueba cae dentro del preentrenamiento \cite{benhenda2026lookahead}; por eso nuestra ventana fuera de muestra empieza tras el corte de entrenamiento del modelo.

El agente que estudiamos es una réplica de AI Hedge Fund \cite{singh2024aihedgefund}, un sistema de cinco personalidades que emulan a inversores reconocidos (Buffett, Wood, Druckenmiller, Burry y Ackman) sobre un mismo modelo de lenguaje. Para cada activo y día cada personalidad emite una señal direccional con un grado de confianza; un agente de riesgo fija el límite de posición según la volatilidad, y un gestor integra las cinco señales y ese límite en una sola decisión, dirección y tamaño, por criterio propio y no por voto. Esa decisión diaria es lo que STRATA supervisa. Sin supervisión el agente reproduce los problemas anteriores: pierde dirección y dinero sobre SPY y la mayoría del panel (Sección~\ref{sec:intro-motivacion}). Lo llamamos M5, y la estructura de su pérdida importa más que el tamaño: mantiene una posición casi constante que apenas atiende a la señal del día. Un fallo sistemático se corrige desde fuera, y ahí entra una capa que reintroduzca esa información.

\section{Supervisión en tiempo de ejecución}
\label{sec:supervision}

La industria ya pide vigilar a estos agentes mientras operan. El discurso profesional lo resume en una consigna, \emph{guardrails, but no handcuffs}: barreras que acoten al agente sin maniatarlo, con monitorización continua, trazas de auditoría y salidas que escalen a un humano \cite{milligan2024guardrails}. Entre los controles aparece, casi de pasada, la \emph{validación estadística}: que las salidas se alineen con normas estadísticas realistas \cite{milligan2024guardrails}. La demanda queda en eslogan, sin decir qué norma, qué contraste ni cómo se mide que la corrección funcione. STRATA lo convierte en detectores falsables con contraste pareado.

Vigilar mientras se decide enlaza con la detección de anomalías \cite{chandola2009anomaly} y con el control en tiempo de ejecución (\emph{runtime enforcement}), que intercepta cada acción y bloquea la que incumple una condición \cite{wang2025agentspec, flehmig2025reliability}. Una regla declarativa codifica qué está prohibido; nosotros modelamos cómo de coherente es una posición con el estado del mercado, con tres detectores estadísticos clásicos e interpretables: RAM mira la coherencia con el régimen (modelo de Markov oculto), PSA la coherencia temporal del agente (puntos de cambio bayesianos) y GSO la compatibilidad del tamaño con la volatilidad (GARCH). Los tres descansan en cómo se describe hoy el régimen de mercado y su riesgo.

\section{Régimen de mercado y riesgo}
\label{sec:regimen}

La idea de fondo de RAM viene de Hamilton: una serie financiera no responde a un único mecanismo, sino que atraviesa estados latentes con propiedades estadísticas distintas y cambia de estado sin que lo veamos, de modo que la distribución de cada retorno depende del estado vigente \cite{hamilton1989}. Condicionar la cartera al régimen mejora las decisiones de forma apreciable \cite{ang2002regime, guidolin2007regime}. Sobre el eje continuo de la volatilidad opera la idea que ancla a GSO: dimensionar la posición para alcanzar una volatilidad objetivo (el \emph{volatility targeting}) mejora las propiedades de la serie de retornos \cite{moreira2017}. La literatura usa de dos \cite{ang2002regime} a cuatro \cite{guidolin2007regime} estados; usamos tres (calma, estrés y crisis): la verosimilitud fuera de muestra descarta dos y por encima de tres se pierde lectura. El régimen funciona como proxy direccional solo de forma condicional, y esa condición sostiene y limita la tesis: el \emph{leverage effect} de los índices agregados hace que en activos como SPY los estados de alta volatilidad coincidan con tramos bajistas, de modo que conocer el régimen informa del signo probable del retorno \cite{black1976, christie1982}. El régimen no predice el retorno: ordena las decisiones del agente. La coincidencia se diluye en valores con apalancamiento débil, donde el régimen pierde poder direccional, y ahí está el límite de la técnica. Estas tres líneas se han tocado por separado; STRATA las une con detectores estadísticos clásicos y exógenos, que el capítulo siguiente formaliza.
